% ============================================
% Title Options
% ============================================

\title{Interspecies Interactions Drive Community-Level Selection in Microbial Coalescence}

% Alternative title:
% Interaction strength determines when microbial communities behave as units of selection

% ============================================
% Authors and Affiliations
% ============================================
\author[1]{Jinyeop Song}
\author[1]{Jiliang Hu}
\author[1]{Jeff Gore}

\affil[1]{Department of Physics, Massachusetts Institute of Technology, Cambridge, MA, USA}

\date{}

\maketitle

% ============================================
% Abstract
% ============================================
\begin{abstract}
It has long been debated in ecology whether communities behave as cohesive units or loose collections of independent species. Here, we study this question through community coalescence---the mixing of previously isolated communities that were each assembled to equilibrium separately. Using bacterial microcosm experiments and ecological modeling, we show that interspecific interaction strength determines whether community-level selection occurs. At moderate to strong interactions, community-level selection is frequent, with one parent community outcompeting the other. At weak interactions, coalescence typically leads to symmetric mixtures of the two communities, indicating that community-level selection disappears. Experiments tuning interaction strength via nutrient concentration confirm these predictions. We further identify two regimes underlying community-level selection: an emergent regime at intermediate strength where collective dynamics shape outcomes in ways not predictable from species traits alone, and a species-driven regime at strong interactions where dominant species determine the winner. Our findings reconcile conflicting observations by showing that interaction strength governs when communities behave as units of selection.
\end{abstract}

\newpage
